\section{The Temptation of the Absolute}

\begin{quotex}
I loved her, and sought her out from my youth, I desired to make her my spouse, and I was a lover of her beauty.
\flright{\textsc{Wisdom 8:2}}

A process of mind is thinking only so far as it realizes the end of thinking, which is always to understand, that is, to see things within a system which renders them necessary.
\flright{\textsc{Brand Blanshard}, \emph{The Nature of Thought}}

The evolution of knowledge is, as Plato long ago portrayed it, the emancipation of individual minds from their accidental limitations, and their education into the knowledge of the one real and intelligible world.
\flright{\textsc{Bernard Bosanquet}, \emph{Logic}}

Truth is free development of one's inner powers. The light is in us.
\flright{\textsc{Benedetto Croce}, \emph{Logic}}
\end{quotex}


In \emph{La Vita Nuova}, \textbf{Dante} tells the story about encountering Beatrice in a public. Not wishing to stare at her beauty in public, he directed his gaze to another woman who was in the same line of sight. That way onlookers would assume he was staring at the other woman even though he could see Beatrice out of the corner of his eye. The woman and the onlookers assumed Dante was interested in the other woman. Not ready to reveal his secret desire, Dante allowed that misconception to persist for several years.

In \emph{Il Convivio} (\emph{The Banquet}), Dante explained that this woman was Philosophy, which was the first dish at the banquet. Ultimately, however, it was Beatrice who led him to Heaven, and not Philosophy.

\paragraph{The Quest for the Absolute}


I, too, understand the attraction of that woman. I have always been interested in Philosophy, not Naturalism which is her ugly step-sister, but the true Philosophy for which I'll use the umbrella term ``Absolute Idealism''. I've been more interested in Plotinus, Spinoza, Schopenhauer, Bradley, Gentile, Collingwood, Blanshard, and more recently, Timothy Sprigge, although there are certainly other options.

In my own case, I would include Indian philosophers like Aurobindo and Sarvepalli Radhakrishnan; the latter note the many points of contact between Eastern philosophy and Western Absolute Idealism.

Fundamentally for Absolute Idealism, thought is primary. The world is intelligible, and there is a drive to know, i.e., to become intelligent. Ultimately, the thinker attains to knowledge of the Absolute in which all antinomies are dissolved. The world, then, in all its manifestations is an appearance of the Absolute.

\paragraph{Emergent Evolution and Holism}


Two books appeared early last century: \emph{Holism and Evolution} by Jan Smuts and \emph{Emergent Evolution} by Lloyd Morgan. Although I doubt they are read much anymore, the two concepts they described are part of today's intellectual atmosphere. Put together, the claim is that evolution proceeds by simpler things organizing into larger wholes (holism), in which new properties emerge, which are inexplicable simply in terms of the constituent elements.

In one sense, these ideas are an advance because emergence is opposed to reductionism and holism to atomism. The atomic viewpoint cannot distinguish between a bunch of matches randomly thrown on a table and a matchbook sculpture\footnote{\url{https://www.bing.com/images/search?q=matchstick+sculpture\&FORM=HDRSC2}}. The sculpture is explained as the result of natural processes and any design is a human superimposition.

However, the realist viewpoint of those two books is the opposite of idealism. They agree that the world consists of ever larger wholes, and that larger whole cannot be reduced to smaller ones. The evolutionary view presumes that there is an inner force of some sort that drives the lower elements to self-organize into the higher. Nevertheless, it cannot be intelligent --- since it hadn't yet evolved --- so the theory is unsatisfactory.

The alternative is that evolution occurs in thought. The intellectual drive is to see things in a larger whole. If the idea is logically coherent, the data of experience are understood at a more satisfactory level. For example, in human affairs, the actions of individuals are more adequately explained in terms of higher relations such as ethnicities, races, political and religious affiliations. To go a step further, time needs to be taken into consideration. Generational changes need to be understood.

On a personal level, the same applies. We note that we are inconsistent in our consciousness. Hence, we should try to develop a unified Self. The more that thinking can encompass, the closer we get to the knowledge of the Absolute.

\paragraph{Pansychism}


This is the idea that interiority is ubiquitous in the world.
\textbf{Rene Guenon} wrote in the \emph{Reign of Quantity}:


\begin{quotex}
there can in fact be no `inanimate' objects in existence, and also that `life' is one of the conditions to which all corporeal existence without exception is subject; and that is why nobody has ever arrived at a satisfactory definition of the difference between the `living' and the `non-living', for that question, like so many others in modern philosophy and science, is only insoluble because there is no good reason for posing it, since the `nonliving' has no place in the domain to which the question is related, and the only differences involved are really no more than mere differences of degree.
\end{quotex}


Obviously, from this perspective, there is no mystery about the origin of life. What science cannot explain is the inner experiences that accompanies life. For science, which deals with quantities, thunder is a sound wave of certain frequencies and energies and lightning is a measurable electromagnetic event. The awareness of thunder and lightning is at best a useless epiphenomenon. Robots will be able to deal with thunder and lightning but will have no inner awareness of it.

A less jaded age would have experienced thunder and lightning as signs of higher intelligences. They believed that natural things like wind, fire, and even bodies of water are the physical manifestations of those higher intelligences.

\paragraph{The Absolute and God}


Croce made the point that Religion and Philosophy were concerned with the same goal: the Absolute. Hence, they are incompatible and once a man has achieved philosophical knowledge, he no longer has any further need of Religion.

Guenon can be misinterpreted that way also, given his claim that all validly Traditional religions, with their symbolism, creeds, and rites, share a common metaphysic. Nevertheless, religion itself cannot be abandoned, since it serves both as a support to metaphysics and a pathway to self-realization.

Many of the Idealists would agree with Croce. Once thought has attained the Absolute, it sees religion as simply an appearance of the Absolute. However, if you accept a version of panpsychism, then the Absolute is not simply an object of thought. It, too, must be conscious and have a Mind, i.e., it is also a subject.

Thinking, then, does not suffice. There must also be a path of realization. Specifically, beyond thinking, a change of being is also required. That leads us beyond philosophy proper.

The Russian theologian \textbf{Sergei Bulgakov} made the transition from Marxism through Idealism to Religion, so he understands what is at stake. In \emph{Sophia: The Wisdom of God}, he explains:

\begin{quotex}
Absolute being, self-existent and self-sufficing, while maintaining all its absolute character, yet establishes as it were alongside or outside of itself a state or relative being, to which it stands as God. The Absolute is God, but God is not the Absolute insofar as the world relates to him.
\end{quotex}



\begin{center}* * *\end{center}

\begin{footnotesize}
\begin{sffamily}
\texttt{james on 2018-02-16 at 04:51 said:}

Thanks for another great post.

\hfill

\texttt{sindre on 2018-02-18 at 11:56 said:}

Christians can learn a lot from other religions; but those of no particular religion cannot learn anything substantive about Christianity -- or indeed anything else.

Instead, these `spiritual seekers' or adherents of the `perennial philosophy' merely become trapped in a self-congratulatory/ self-indulgent version of the mainstream modern lifestyle -- with its trapping of the sexual revolution, political Leftism, and technological self-manipulation of emotions (via drink, drugs, `body-art', social activism etc.).

The perennial philosophy is therefore an amusement, not a conviction; an open-ended life-option that seeks no end and attains no progression; an interest, not a faith; and complacency, not courage. It is simply an addictive kind of feebleness -- a craving for a pain-free and engaging mortal life -- with genuine escape/ enlightenment always just around the next corner... True seeking is based on the conviction that there is a right answer that leads to more right answers -- it is serial finding, not serial seeking.

All answers -- all true knowledge -- are partial and biased (i.e. perspectival). Progress comes from the process of of integrating and correcting these answers -- and it never ends; because reality is creative, so there is always more to know.

Spiritual seekers expend their mortal lives in looking-for a form -- asserting that all forms and motivations are one. This is lethal to real spirituality; since it discards the ens and essences of metaphysics and motivation -- and instead focuses on the means and peripheries of emotions, lifestyle, and detached/ fragmentary utterances and writings.

(A monk who seeks detachment-from the world is equated with a shaman who seeks absorption-into the abstract divine and a priest who seeks a personal communion-with a personal god... on the basis that they all wear robes and spend time sitting with their eyes closed... )

That the perennial philosophy is nonsense is a matter of simple common sense and logic; which is why the idea is only ever held by intellectuals who have the cognitive capacity to confuse themselves with complexity.

The perennial philosophy is at least tolerated, and probably encouraged, by the demonic beings who pursue strategic evil in this world, because it permanently neutralises the basic human quest for god, for true religion. Once enmeshed in endless syncretism and research -- the spiritual seeker is trapped in a kaleidoscope of images and assertions, and cannot find a way out.

Of course, even sincere and genuinely motivated individuals will likely go through some kind of phase of comparison and learning about religions -- but to defend this as a valid permanent option, to crystallise it as `the truth' about the human condition, is fatal to faith, hope and charity -- leading instead to this-worldly materialism, purposive hedonism and despair.

\end{sffamily}
\end{footnotesize}

